%%=============================================================================
%% Samenvatting
%%=============================================================================

% TODO: De "abstract" of samenvatting is een kernachtige (~ 1 blz. voor een
% thesis) synthese van het document.
%
% Een goede abstract biedt een kernachtig antwoord op volgende vragen:
%
% 1. Waarover gaat de bachelorproef?
% 2. Waarom heb je er over geschreven?
% 3. Hoe heb je het onderzoek uitgevoerd?
% 4. Wat waren de resultaten? Wat blijkt uit je onderzoek?
% 5. Wat betekenen je resultaten? Wat is de relevantie voor het werkveld?
%
% Daarom bestaat een abstract uit volgende componenten:
%
% - inleiding + kaderen thema
% - probleemstelling
% - (centrale) onderzoeksvraag
% - onderzoeksdoelstelling
% - methodologie
% - resultaten (beperk tot de belangrijkste, relevant voor de onderzoeksvraag)
% - conclusies, aanbevelingen, beperkingen
%
% LET OP! Een samenvatting is GEEN voorwoord!

%%---------- Nederlandse samenvatting -----------------------------------------
%
% TODO: Als je je bachelorproef in het Engels schrijft, moet je eerst een
% Nederlandse samenvatting invoegen. Haal daarvoor onderstaande code uit
% commentaar.
% Wie zijn bachelorproef in het Nederlands schrijft, kan dit negeren, de inhoud
% wordt niet in het document ingevoegd.

\IfLanguageName{english}{%
\selectlanguage{dutch}
\chapter*{Samenvatting}
\lipsum[1-4]
\selectlanguage{english}
}{}

%%---------- Samenvatting -----------------------------------------------------
% De samenvatting in de hoofdtaal van het document

\chapter*{\IfLanguageName{dutch}{Samenvatting}{Abstract}}

Mainframes zijn vandaag nog steeds een cruciaal onderdeel van de bedrijfskritieke infrastructuur binnen verschillende sectoren, 
variërend van financiële bedrijven zoals Euroclear tot de staalindustrie met ArcelorMittal.
Toch kampt de mainframe industrie als geheel met een grote vergrijzings crisis.
De laatse jaren neemt het aantal ontwikkelaars met kennis van PL/1 en CICS sterk af door vergrijzing en dalende instroom van nieuw talent uit het onderwijs.

Deze bachelorproef bekijkt hoe dit structureel probleem kan aangepakt worden voor PL/1 en CICS ontwikkelaars op z/OS.
Het hoofddoel van deze bachelorproef is het ontwikkelen van een competentieframework, dat een overzicht geeft van kennis en vaardigheden die een startende PL/1 en CICS developer nodig heeft.
Daarnaast ook het ontwikkelen van een proof of concept leerplatform dat studenten in staat moet stellen om deze competenties stapsgewijs op te bouwen.

Het framework werd opgesteld op basis van een literatuurstudie en semigestructureerde interviews met professionals uit Euroclear.
De proof of concept is een online leerplatform dat is opgebouwd uit verschillende modules, gebaseerd op het opgestelde competentieframework.
Binnen elke module verwerken studenten de leerstof aan de hand van theorie, praktische programmeeroefeningen en quizzen.


Ten laatste werd deze leersite samen met het competentie framework gevalideerd door experten uit het werkveld en mede studenten van Hogent.

De resultaten van deze validatie tonen aan dat een gericht framework voor PL/1 en CICS ontwikkelaar een reële meerwaarde brengen voor zowel onderwijsinstellingen als bedrijven die nieuw talent willen aanwerven in hun legacy systemen.
Door een praktijkgerichte aanpak te hanteren via een leerplatform met quizzen en oefeningen, wordt de leerstof niet alleen effectiever overgebracht, maar resulteert dit vooral in een veel betere kennisretentie bij de studenten.
Het verkleint de skill gap en maakt het leren van deze complexe technologien meer toegankelijk.
Hierdoor krijgen startende ontwikkelaars een duidelijk groeipad om snel productief te worden in hun rol.

